\documentclass[11pt]{article}

% ---------------------------------------------------------------------------
% AI Applications in Research: Concepts from Scott Cunningham's Claude Code series
% Compiled as a conceptual summary of agentic AI for empirical research.
% ---------------------------------------------------------------------------

\usepackage[utf8]{inputenc}
\usepackage[T1]{fontenc}
\usepackage[margin=1in]{geometry}
\usepackage{amsmath,amssymb}
\usepackage{enumitem}
\usepackage{booktabs}
\usepackage{xcolor}
\usepackage[colorlinks=true,linkcolor=blue,urlcolor=blue,citecolor=blue]{hyperref}
\usepackage{titlesec}
\usepackage{parskip}

\titleformat{\section}{\large\bfseries}{\thesection}{1em}{}
\titleformat{\subsection}{\normalsize\bfseries}{\thesubsection}{1em}{}

\title{\bfseries AI Applications in Research:\\
\large Agentic Coding and the Case of Claude Code}
\author{Summary compiled from Scott Cunningham's \emph{Scott's Mixtape Substack}}
\date{August 10, 2026}

\begin{document}
\maketitle

\begin{abstract}
\noindent This document summarizes the main concepts surrounding the use of
artificial intelligence---and in particular \emph{agentic} coding tools such as
Claude Code---in empirical and quantitative research. It draws on the ongoing
essay series by economist Scott Cunningham, who has documented (in 60+ essays)
how coding agents reshape the production, verification, and understanding of
research. The central thesis is that AI agents are \emph{not} chatbots: they are
a distinct ``harness'' that can autonomously execute an end-to-end research
pipeline, which shifts the binding constraint in research from production to
verification and understanding.
\end{abstract}

\section{Three Harnesses for Language Models}
A useful framing distinguishes three ways language models are packaged for users,
each a different \emph{harness} around the same underlying model:
\begin{enumerate}[label=\arabic*.]
  \item \textbf{Chatbots} --- e.g.\ ChatGPT and Claude. Conversational, text-in/text-out.
  \item \textbf{Coding copilots} --- e.g.\ Cursor and GitHub Copilot. Inline code suggestion inside an editor.
  \item \textbf{Coding agents} --- e.g.\ Claude Code and Codex. Autonomous tools that read and write files, run shell commands, download data, and orchestrate multi-step work.
\end{enumerate}
The key conceptual lesson---and one of the most valuable takeaways for a new
audience---is that \textbf{AI agents for research are not chatbots}. Prior
familiarity with ChatGPT can \emph{hinder} understanding, because it leads people
to believe the goal is finding the ``right prompt.'' In practice, one does not
``prompt-engineer'' an agent so much as \emph{converse} with it as a thinking
partner; ``prompt engineering,'' if it ever existed, is not how agentic research
is done.

\section{What a Coding Agent Actually Does}
Unlike a chatbot, a coding agent operates directly on your computer. Conceptually,
it can:
\begin{itemize}[leftmargin=1.5em]
  \item Access the entire hard drive, files, and tools of the machine.
  \item Execute a wide range of shell/bash commands autonomously.
  \item Download datasets, run analyses across multiple languages and packages, and replicate pipelines from scratch.
  \item Produce a full manuscript (e.g.\ a 20-page double-spaced draft) with tables, figures, references, robustness checks in appendices, and an accompanying slide deck.
\end{itemize}
Cunningham illustrates this with a live ``magic trick'': starting from an
\emph{empty folder} (``nothing up my sleeve''), a single prompt directs the agent
to articulate a research question and identification strategy, gather data, run
the analysis, and write up the results---all in 20--30 minutes---while the
presenter lectures in parallel.

\section{The Three Pillars of Research with Agents}
Once production becomes nearly costless, the intellectual work relocates.
Cunningham frames modern research around three tasks:
\begin{center}
\begin{tabular}{@{}ll@{}}
\toprule
\textbf{Pillar} & \textbf{Question it answers} \\
\midrule
Producing research   & Can the pipeline be generated and run? \\
Verifying research   & Are the results correct---and what does ``correct'' mean here? \\
Understanding research & Do \emph{I} understand what was done and why? \\
\bottomrule
\end{tabular}
\end{center}
The magic trick makes the stakes visceral: the audience sees figures and tables
that \emph{exist} without anyone knowing how they were made or whether they are
accurate. When human time inputs fall to zero, \textbf{verification and
understanding---not production---become the binding constraints}. ``Understanding
is the new bottleneck.''

\section{An Economic Model of Research Production}
Cunningham situates agents within standard microeconomic tools rather than
treating hooks and skills as isolated tricks:
\begin{itemize}[leftmargin=1.5em]
  \item \textbf{Isoquants and inputs.} Research output is produced from \emph{machine time} and \emph{human time}. Agents dramatically raise the productivity of---and substitute for---human time, shifting the isoquant and lowering the marginal cost of empirical work.
  \item \textbf{Stages of research.} He invokes Shockley's classic analysis of the productivity of scientists (the 7--8 stages of research) to locate where agents intervene.
  \item \textbf{Supporting concepts.} Demand curves, surplus, and flow charts help visualize how the technology reallocates effort and value across the research process.
\end{itemize}

\section{Workflow: Continuous Verification and Understanding}
Because correctness is no longer guaranteed by the labor of production, disciplined
workflow becomes essential. Recurring elements in the series include:
\begin{itemize}[leftmargin=1.5em]
  \item \textbf{Markdown / \texttt{CLAUDE.md} instructions} --- persistent, project-level context and conventions for the agent.
  \item \textbf{Hooks} --- automated actions the harness runs at defined points (e.g.\ before/after a step) to enforce checks.
  \item \textbf{Skills} --- packaged, reusable procedures the agent can invoke for recurring tasks.
  \item \textbf{A checklist-based dashboard / GUI} --- to manage the workflow and enforce \emph{continuous understanding} and \emph{continuous verification} rather than a one-time review at the end.
  \item \textbf{Guarding against stale code and deprecation} --- structured workflows matter when, for example, an analytical sample silently changes size ($N \to N+k$) across a long project.
  \item \textbf{Coding audits} --- independently re-implementing a pipeline in multiple languages and packages to cross-check results.
\end{itemize}

\section{Pedagogy: How to Teach Agentic Research}
Cunningham's teaching template ``defamiliarizes the familiar'' and brings the
median audience member forward, not just the few:
\begin{enumerate}[label=\arabic*.]
  \item \textbf{History of thought \& narrative} --- e.g.\ the origin story of Claude Code (created at Anthropic in late 2024 by Boris Cherny for internal use; public via the CLI in 2025; adoption by empiricists accelerating once it reached the desktop app in late 2025, since social scientists rarely used the terminal).
  \item \textbf{Numerical examples \& interactive demos} --- the live ``magic trick'' from an empty folder, which commands undivided attention.
  \item \textbf{Video / interactivity} --- to sustain engagement.
  \item \textbf{Technical rigor last} --- installation, the underlying bash commands, and a conceptual framework for markdowns, hooks, and skills, positioned inside a model of what research \emph{is}.
\end{enumerate}
The goals on the ground are emotional and cognitive at once: emotion, insight,
understanding, feeling personally seen and validated, and engagement---reached by
\emph{backwards induction} from those outcomes through the five teaching principles.

\section{Open Questions and Debates}
\begin{itemize}[leftmargin=1.5em]
  \item \textbf{Prompt disclosure.} Requiring authors to ``submit their prompts'' with a paper is argued to be misguided: with agents used as thinking partners, transcripts are more like private conversations with coauthors than a reproducible artifact, and mandatory disclosure raises privacy and adverse-selection concerns.
  \item \textbf{Learning by doing.} There is ``no such thing as prompt engineering'' for agents; the relevant skill is natural, emotionally intelligent conversation, largely acquired through doing real projects.
  \item \textbf{The verification frontier.} As production cost falls, the discipline's open problem is building reliable, continuous methods to verify and truly understand agent-produced results.
\end{itemize}

\section{Key Takeaways}
\begin{enumerate}[label=\textbf{\arabic*.}]
  \item AI agents are a distinct harness, categorically different from chatbots.
  \item Coding agents can autonomously execute an end-to-end empirical pipeline from an empty folder.
  \item The marginal cost of producing empirical research has collapsed; the new bottlenecks are \emph{verification} and \emph{understanding}.
  \item Research now rests on three pillars: producing, verifying, and understanding.
  \item Standard micro theory (isoquants, human/machine time, stages of research) offers a rigorous frame for the shift.
  \item Disciplined workflow---\texttt{CLAUDE.md}, hooks, skills, and checklist dashboards---delivers continuous verification and understanding.
  \item Effective teaching leads with history and narrative and a live demonstration before technical detail.
\end{enumerate}

\section*{Source}
\addcontentsline{toc}{section}{Source}
Scott Cunningham, ``My strategy for talks on agents for research,''
\emph{Scott's Mixtape Substack}, August 10, 2026, and the broader Claude Code
essay series. \url{https://causalinf.substack.com/p/my-strategy-for-talks-on-agents-for}

\end{document}
